\chapter{Esercizi: Collaudo, tipi di dato e input in Java}
\label{chap:esercizi-collaudo-tipi-input-java}

\begin{exercise}
\label{ex:fi-esercizi-java-collaudo-passi}
Quali sono i passi principali per costruire un programma di collaudo di una
classe?
\end{exercise}

Si definisce una classe di collaudo con un metodo \texttt{main}, si creano gli
oggetti necessari e se ne invocano i metodi con ingressi scelti. I risultati
ottenuti vengono confrontati con quelli attesi, considerando anche i casi
limite. Il superamento dei test non dimostra la correttezza generale
(\cref{sec:fi-java-collaudo}).

\begin{exercise}
\label{ex:fi-esercizi-java-collaudo-package}
Che cosa sono i \emph{package} in Java e a che cosa servono?
\end{exercise}

Sono raggruppamenti di classi correlate, utilizzati per organizzare il codice
e distinguere classi che possono avere lo stesso nome semplice. Per esempio,
\texttt{Rectangle} appartiene al pacchetto \texttt{java.awt}
(\cref{sec:fi-java-collaudo}).

\begin{exercise}
\label{ex:fi-esercizi-java-collaudo-import}
Scrivi la sintassi generale di un'istruzione \texttt{import}.
\end{exercise}

La forma generale è \texttt{import nomePacchetto.NomeClasse;}.

Per esempio, \texttt{import java.awt.Rectangle;} consente di utilizzare
il nome semplice della classe senza ripeterne il nome completo del pacchetto
(\cref{sec:fi-java-collaudo}).

\begin{exercise}
\label{ex:fi-esercizi-java-collaudo-java-lang}
Perché le classi \texttt{System} e \texttt{String} possono essere usate senza
un'importazione esplicita?
\end{exercise}

Appartengono al pacchetto \texttt{java.lang}, le cui classi sono disponibili
automaticamente mediante il loro nome semplice
(\cref{sec:fi-java-collaudo}).

\begin{exercise}
\label{ex:fi-esercizi-java-collaudo-primitivi-oggetti}
Qual è la differenza tra tipi di dato fondamentali e oggetti in Java?
\end{exercise}

Una variabile di tipo primitivo contiene direttamente un valore del tipo,
mentre una variabile di tipo riferimento contiene un riferimento a un oggetto
o \texttt{null}. Gli oggetti sono istanze di classi e possono esporre metodi
per operare sul proprio stato
(\cref{sec:fi-java-tipi-fondamentali,sec:fi-java-riferimenti}).

\begin{exercise}
\label{ex:fi-esercizi-java-collaudo-tipi-numerici}
Elenca i principali tipi di dato fondamentali numerici di Java.
\end{exercise}

I tipi interi con segno sono \texttt{byte}, \texttt{short}, \texttt{int} e
\texttt{long}, rispettivamente di 8, 16, 32 e 64 bit. I tipi in virgola
mobile sono \texttt{float} e \texttt{double}, rispettivamente di 32 e 64 bit
(\cref{tab:fi-java-tipi-primitivi}).

\begin{exercise}
\label{ex:fi-esercizi-java-collaudo-overflow-int}
Che cosa succede in Java quando una somma tra \texttt{int} provoca
\emph{overflow}?
\end{exercise}

Il risultato subisce un \emph{wrap-around} nell'intervallo degli interi a
32 bit, senza generare un'eccezione. Per esempio, sommando \(1\) al massimo valore di \texttt{int} si ottiene
il minimo valore dello stesso tipo
(\cref{sec:fi-java-overflow,sec:fi-complemento-due}).

\begin{exercise}
\label{ex:fi-esercizi-java-collaudo-metodi-exact}
Quali metodi della classe \texttt{Math} si possono usare per intercettare
l'overflow aritmetico?
\end{exercise}

I metodi \texttt{Math.addExact()} e \texttt{Math.multiplyExact()} eseguono
rispettivamente addizione e moltiplicazione di interi, sollevando
\texttt{ArithmeticException} se il risultato non è rappresentabile
(\cref{sec:fi-java-overflow}).

\begin{exercise}
\label{ex:fi-esercizi-java-collaudo-arithmetic-exception}
Quali operazioni aritmetiche tra interi possono generare una
\texttt{ArithmeticException}?
\end{exercise}

Fra gli operatori aritmetici ordinari, sia la divisione intera
\texttt{/} sia il resto \texttt{\%} generano
\texttt{ArithmeticException} quando il divisore è zero.
Il normale overflow di somma e prodotto non genera eccezioni;
possono invece generarle metodi come \texttt{Math.addExact()}
e \texttt{Math.multiplyExact()}
(\cref{sec:fi-java-overflow}).

\begin{exercise}
\label{ex:fi-esercizi-java-collaudo-letterale-long}
Perché per assegnare un valore molto grande a una variabile \texttt{long}
bisogna aggiungere la lettera \texttt{L}?
\end{exercise}

Un letterale intero decimale senza suffisso è di tipo \texttt{int}; se
supera l'intervallo di tale tipo, non è valido come letterale
\texttt{int}, anche quando la variabile di destinazione è
\texttt{long}. Il suffisso \texttt{L}, come in
\texttt{3000000000L}, indica un letterale \texttt{long}
(\cref{sec:fi-java-tipi-fondamentali}).

\begin{exercise}
\label{ex:fi-esercizi-java-collaudo-float-double}
Qual è la differenza tra \texttt{float} e \texttt{double}?
\end{exercise}

\texttt{float} usa il formato IEEE 754 \emph{binary32}, con 32 bit e circa
7 cifre decimali significative; \texttt{double} usa \emph{binary64}, con
64 bit e circa 16 cifre. Quest'ultimo offre maggiore precisione e un
intervallo di valori più ampio
(\cref{sec:fi-java-precisione,sec:fi-virgola-mobile}).

\begin{exercise}
\label{ex:fi-esercizi-java-collaudo-letterale-float}
Perché per assegnare un valore decimale a una variabile \texttt{float}
bisogna aggiungere la lettera \texttt{f}?
\end{exercise}

Un letterale decimale senza suffisso è di tipo \texttt{double}, la cui
conversione a \texttt{float} richiede generalmente un cast esplicito.
Il suffisso \texttt{f} specifica direttamente un letterale
\texttt{float}, come in \texttt{float x = 30.2f;}
(\cref{sec:fi-java-precisione,sec:fi-java-conversioni}).

\begin{exercise}
\label{ex:fi-esercizi-java-collaudo-nan-infinita}
Che cosa rappresentano \texttt{NaN}, \texttt{POSITIVE\_INFINITY} e
\texttt{NEGATIVE\_INFINITY}?
\end{exercise}

\texttt{NaN} rappresenta un risultato non numerico, come
\texttt{0.0 / 0.0}; le altre due costanti rappresentano le
infinità positiva e negativa, che possono derivare, per esempio,
da una divisione in virgola mobile per zero. Sono valori speciali
dei tipi \texttt{float} e \texttt{double}, disponibili come costanti delle classi \texttt{Float} e
\texttt{Double} (\cref{sec:fi-java-overflow,tab:fi-valori-speciali}).

\begin{exercise}
\label{ex:fi-esercizi-java-collaudo-cast-435}
Spiega perché \texttt{(int)(100 * 4.35)} può dare \texttt{434}
invece di \texttt{435}.
\end{exercise}

Il valore \(4.35\) è approssimato in \texttt{double}; il prodotto può
valere \texttt{434.99999999999994}. Il cast a \texttt{int} tronca verso zero,
restituendo \texttt{434}, anziché arrotondare al più vicino
(\cref{sec:fi-java-precisione,sec:fi-java-conversioni}).

\begin{exercise}
\label{ex:fi-esercizi-java-collaudo-conversione-cast}
Che differenza c'è tra una conversione implicita e un \emph{cast}?
\end{exercise}

Una conversione implicita è applicata automaticamente quando è ammessa
dalle regole del linguaggio, per esempio da \texttt{int} a
\texttt{long}. Un cast richiede di indicare esplicitamente il tipo di
destinazione, come in \texttt{(int) x}; può comportare perdita di
informazione (\cref{sec:fi-java-conversioni}).

\begin{exercise}
\label{ex:fi-esercizi-java-collaudo-cast-int}
Che cosa fa il cast \texttt{(int)} applicato a un numero in virgola mobile?
\end{exercise}

Per un numero finito nell'intervallo rappresentabile da
\texttt{int}, elimina la parte frazionaria troncando verso zero:
\texttt{(int) 2.95} vale \texttt{2} e
\texttt{(int) (-2.95)} vale \texttt{-2}. Non effettua un
arrotondamento all'intero più vicino
(\cref{sec:fi-java-conversioni}).

\begin{exercise}
\label{ex:fi-esercizi-java-collaudo-math-round}
Perché \texttt{Math.round()} è preferibile rispetto a
\texttt{(int)(x + 0.5)} per arrotondare?
\end{exercise}

\texttt{Math.round()} esprime direttamente l'arrotondamento
all'intero più vicino, con scelta dell'intero maggiore in caso
di equidistanza. La formula con il cast non funziona in generale
per i negativi: per \(x=-2.95\) restituisce \(-2\), mentre
\texttt{Math.round(x)} restituisce \(-3\)
(\cref{sec:fi-java-conversioni}).

\begin{exercise}
\label{ex:fi-esercizi-java-collaudo-big-number}
A che cosa servono le classi \texttt{BigInteger} e
\texttt{BigDecimal}?
\end{exercise}

\texttt{BigInteger} rappresenta interi di dimensione non limitata
dalla larghezza di \texttt{long}; \texttt{BigDecimal} consente
l'aritmetica decimale a precisione arbitraria e il controllo
dell'arrotondamento. Entrambe appartengono a \texttt{java.math}
e utilizzano metodi per le operazioni
(\cref{sec:fi-java-conversioni}).

\begin{exercise}
\label{ex:fi-esercizi-java-collaudo-costante}
Scrivi la sintassi generale per dichiarare una costante in Java.
\end{exercise}

Si utilizza \texttt{final tipo NOME\_COSTANTE = espressione;},
per esempio \texttt{final double PI\_GRECO = 3.14;}.
La variabile può essere assegnata una sola volta; se contiene un
riferimento, \texttt{final} non rende immutabile l'oggetto
(\cref{sec:fi-java-costanti-aritmetica}).

\begin{exercise}
\label{ex:fi-esercizi-java-collaudo-divisione-resto}
Qual è la differenza tra divisione intera \texttt{/} e operatore
resto \texttt{\%}?
\end{exercise}

Con operandi interi, \texttt{/} restituisce il quoziente troncato
verso zero e \texttt{\%} il resto. Per esempio,
\texttt{235 / 100} vale \texttt{2}, mentre
\texttt{235 \% 100} vale \texttt{35}
(\cref{sec:fi-java-costanti-aritmetica}).

\begin{exercise}
\label{ex:fi-esercizi-java-collaudo-math-classe}
A che cosa serve la classe \texttt{Math}?
\end{exercise}

Fornisce metodi statici per operazioni matematiche, tra cui
\texttt{pow()}, \texttt{sqrt()}, \texttt{log()} e \texttt{sin()},
oltre alle costanti \texttt{Math.PI} e \texttt{Math.E}
(\cref{sec:fi-java-math}).

\begin{exercise}
\label{ex:fi-esercizi-java-collaudo-metodi-statici-oggetto}
Qual è la differenza tra un metodo statico come
\texttt{Math.round()} e un metodo invocato su un oggetto come
\texttt{System.out.println()}?
\end{exercise}

Un metodo statico appartiene alla classe ed è invocabile senza
crearne un'istanza, come \texttt{Math.round()}. Invece,
\texttt{println()} è un metodo dell'oggetto cui si riferisce
\texttt{System.out}, un campo statico della classe \texttt{System}
(\cref{sec:fi-java-math,sec:fi-java-cicli-metodi-statici,sec:fi-java-invocazione-metodi}).

\begin{exercise}
\label{ex:fi-esercizi-java-collaudo-substring}
Come funziona il metodo \texttt{substring(startIndex, endIndex)}?
\end{exercise}

Restituisce una nuova stringa formata dalle unità di codice UTF-16
con indici da \texttt{startIndex}, incluso, a \texttt{endIndex},
escluso. Per una stringa di lunghezza \(n\) si richiede
\(0\le\texttt{startIndex}\le\texttt{endIndex}\le n\);
la lunghezza della sottostringa è la differenza dei due indici
(\cref{sec:fi-java-substring}).

\begin{exercise}
\label{ex:fi-esercizi-java-collaudo-string-index-exception}
In quali casi si può ottenere una
\texttt{StringIndexOutOfBoundsException}?
\end{exercise}

Quando si usa un indice non valido in operazioni sulle stringhe:
per \texttt{substring()}, se l'inizio è negativo, la fine supera
la lunghezza oppure l'inizio è maggiore della fine; per
\texttt{charAt(i)}, se \(i<0\) o \(i\ge n\), dove \(n\) è la
lunghezza della stringa
(\cref{sec:fi-java-substring,sec:fi-java-char-escape}).
